% !TEX root = main.tex
% -----------------------------------------------------------------------------
% Chapter 1: Introduction
% -----------------------------------------------------------------------------

\chapter{Introduction}
\label{ch:introduction}

\section{Background}
Environmental monitoring systems now produce streams of readings at a pace that would have been awkward to manage only a few years ago. A single sensor can emit measurements every few seconds or minutes, and a modest deployment quickly becomes a long sequence of temperature, humidity, light, and voltage values that must be stored, aligned, cleaned, and interpreted. That scale is useful, but it is not tidy. Real sensor data are often interrupted by communication dropouts, battery failures, calibration drift, or brief hardware faults, and the result is a series with gaps that may be small, irregular, or unexpectedly long. Missing values are rarely the only issue either; corrupted readings and obvious spikes can sit beside them, turning the observed signal into a mixture of genuine measurements and error rather than a simple incomplete record \cite{Adhikari2022ACS,Alwateer2024,Agbo2022,Osman2018}.

\subsection{Why missingness matters}
The difficulty is not just that values are missing, but that the mechanism behind the missingness may matter. Classical missing-data theory separates missing completely at random, missing at random, and missing not at random, and the same imputation method can behave quite differently across those settings \cite{Little2019,Pham2022,Jamshidian2014}. Sensor data often falls somewhere between these categories. A power outage may remove a long block of readings. A noisy device may continue to report values that look plausible in isolation but become suspicious once the surrounding context is visible. That ambiguity makes the problem less neat than the textbook version, and probably closer to what field engineers and applied researchers actually face.

\subsection{Patterns in time-series gaps}
Time-series imputation is especially sensitive to the pattern of missingness. Filling a one-step gap in the middle of a smooth curve is not the same task as reconstructing the first few samples of a day, where there is no left-hand context at all. The same applies to longer gaps: interpolation may look acceptable when the missing interval is short, yet it can flatten peaks or miss a sharp change when the gap widens. Recent studies on IoT and industrial sensor data suggest that the location, length, and structure of the gap can materially change the performance of both statistical and learned methods \cite{Wu2022,Ahmed2022,Ahmed2022LongGM,Velasco-Gallego2020,Rosero-Montalvo2023}.

\section{Motivation}
This thesis focuses on a univariate temperature series, but that narrower scope does not make the problem trivial. Temperature readings are often smoother than other sensor channels, which may suggest that simple interpolation is enough. Sometimes it is. Yet a smooth daily curve can still hide abrupt changes caused by weather shifts, heater cycles, or sensor drift. A method that blurs everything can produce a plausible line and still miss the structure that matters. That tension is one reason imputation remains a useful research problem rather than a solved preprocessing step.

The literature has moved in several directions at once. Surveys of missing-data imputation describe a broad field that ranges from classical statistical techniques to machine learning, deep learning, and generative models, with no single family dominating every use case \cite{Adnan2022ARO,Miao2023,Alwateer2024,Adhikari2022ACS,Fan2025AnIA}. Simple methods such as mean, median, forward fill, and interpolation remain common because they are easy to explain and inexpensive to compute. Machine-learning methods, including K-nearest neighbours, random forests, gradient boosting, and iterative regression-based imputers, may perform better when the signal has a learnable structure. Deep approaches, especially those built around recurrent networks, autoencoders, or adversarial training, often promise stronger reconstructions on difficult gaps, though that promise depends heavily on the dataset, the masking regime, and the care taken during evaluation \cite{Okafor2021,Yoon2018GAINMD,Wang2024MissingDI,Shahbazian2023GenerativeAN}.

Those trade-offs are not abstract. In IoT and environmental monitoring, an imputation method has to be accurate, stable, explainable enough to justify to a reviewer or supervisor, and affordable enough to run on modest hardware. A method that improves RMSE by a small amount but doubles training time may still be useful in an offline study, yet less attractive if the same pipeline must later run on live data. By the same token, a very simple method may be appealing for deployment and still fail to preserve the local shape of the series. Comparative work on sensor data has repeatedly shown that the best method is often contingent on the data source, the gap pattern, and the evaluation design \cite{Koshta2021,Erhan2021CriticalCO,Lee2025ComparisonOM,Farjallah2025EvaluationOM}.

\section{Research Aim and Objectives}
The aim of this thesis is to examine that balance in a concrete setting. Using a real environmental sensor dataset, the study introduces controlled missingness into a temperature series, treats selected outliers as values to be reconstructed, and compares a set of baseline imputers with a deep sequence model based on an LSTM Variational Autoencoder. The main model is trained with sliding windows and mini-batch gradient descent, which is likely to be more practical than fitting an entire long sequence in one pass when the series grows. The decoder then reconstructs the signal, and overlapping window predictions are stitched together for evaluation. That design is not the only viable choice, but it offers a sensible compromise between computational cost and sequential context.

\subsection{Research objectives}
The work is guided by the following objectives:

\begin{itemize}
	\item prepare a reproducible pipeline for loading, sorting, and masking the sensor data;
	\item compare standard statistical imputers with machine-learning baselines on the same missingness pattern;
	\item implement and train an LSTM-VAE for sequence reconstruction;
	\item evaluate the methods using masked-point error metrics such as MAE and RMSE;
	\item inspect whether the learned model actually preserves local structure better than simpler alternatives.
\end{itemize}

\subsection{Research question}
The central research question is direct: when the same temperature series is corrupted in the same way, does the LSTM-VAE produce a better reconstruction than the baselines, and if so, under what kinds of gaps does that advantage appear most clearly? The answer may not be identical across all missingness patterns. In fact, one useful outcome of the study may be discovering where the deep model does not help as much as expected.

\section{Scope and Contributions}
The scope is deliberately narrow. The work concentrates on a single sensor and a single temperature channel so that the comparison remains interpretable and the effect of each imputation method stays visible. That choice leaves multivariate sensor fusion, irregular sampling, and streaming deployment outside the present study. Those extensions are interesting, but they would blur the central comparison and make the chapter harder to read without strengthening the core argument.

Within that boundary, the thesis makes three modest but useful contributions. First, it sets out a fair experimental protocol in which all methods face the same mask and the same ground truth. Second, it collects statistical, machine-learning, and deep-learning imputers in one place so their differences can be compared without shifting assumptions from one implementation to another. Third, it explores an LSTM-VAE design for time-series imputation that is trained and evaluated in a windowed fashion, which may make the method easier to scale than a single full-sequence pass.

\section{Thesis Structure}
The remainder of the thesis is organised as follows. Chapter 2 reviews related work on missing-data mechanisms, time-series imputation, and deep learning approaches for sensor data. Chapter 3 describes the dataset, masking strategy, baseline methods, and the proposed LSTM-VAE pipeline. Chapter 4 presents the experimental results and compares the methods quantitatively and visually. Chapter 5 discusses what the results suggest, where the methods appear to fail, and how the design might be extended. Chapter 6 concludes the thesis and points to directions for future work.

Time-series imputation is especially sensitive to the pattern of missingness. Filling a one-step gap in the middle of a smooth curve is not the same task as reconstructing the first few samples of a day, where there is no left-hand context at all. The same applies to longer gaps: interpolation may look acceptable when the missing interval is short, yet it can flatten peaks or miss a sharp change when the gap widens. Recent studies on IoT and industrial sensor data suggest that the location, length, and structure of the gap can materially change the performance of both statistical and learned methods \cite{Wu2022,Ahmed2022,Ahmed2022LongGM,Velasco-Gallego2020,Rosero-Montalvo2023}.

This thesis focuses on a univariate temperature series, but that narrower scope does not make the problem trivial. Temperature readings are often smoother than other sensor channels, which may suggest that simple interpolation is enough. Sometimes it is. Yet a smooth daily curve can still hide abrupt changes caused by weather shifts, heater cycles, or sensor drift. A method that blurs everything can produce a plausible line and still miss the structure that matters. That tension is one reason imputation remains a useful research problem rather than a solved preprocessing step.

The literature has moved in several directions at once. Surveys of missing-data imputation describe a broad field that ranges from classical statistical techniques to machine learning, deep learning, and generative models, with no single family dominating every use case \cite{Adnan2022ARO,Miao2023,Alwateer2024,Adhikari2022ACS,Fan2025AnIA}. Simple methods such as mean, median, forward fill, and interpolation remain common because they are easy to explain and inexpensive to compute. Machine-learning methods, including K-nearest neighbours, random forests, gradient boosting, and iterative regression-based imputers, may perform better when the signal has a learnable structure. Deep approaches, especially those built around recurrent networks, autoencoders, or adversarial training, often promise stronger reconstructions on difficult gaps, though that promise depends heavily on the dataset, the masking regime, and the care taken during evaluation \cite{Okafor2021,Yoon2018GAINMD,Wang2024MissingDI,Shahbazian2023GenerativeAN}.

Those trade-offs are not abstract. In IoT and environmental monitoring, an imputation method has to be accurate, stable, explainable enough to justify to a reviewer or supervisor, and affordable enough to run on modest hardware. A method that improves RMSE by a small amount but doubles training time may still be useful in an offline study, yet less attractive if the same pipeline must later run on live data. By the same token, a very simple method may be appealing for deployment and still fail to preserve the local shape of the series. Comparative work on sensor data has repeatedly shown that the best method is often contingent on the data source, the gap pattern, and the evaluation design \cite{Koshta2021,Erhan2021CriticalCO,Lee2025ComparisonOM,Farjallah2025EvaluationOM}.

The aim of this thesis is to examine that balance in a concrete setting. Using a real environmental sensor dataset, the study introduces controlled missingness into a temperature series, treats selected outliers as values to be reconstructed, and compares a set of baseline imputers with a deep sequence model based on an LSTM Variational Autoencoder. The main model is trained with sliding windows and mini-batch gradient descent, which is likely to be more practical than fitting an entire long sequence in one pass when the series grows. The decoder then reconstructs the signal, and overlapping window predictions are stitched together for evaluation. That design is not the only viable choice, but it offers a sensible compromise between computational cost and sequential context.

The work is guided by the following objectives:

\begin{itemize}
	\item prepare a reproducible pipeline for loading, sorting, and masking the sensor data;
	\item compare standard statistical imputers with machine-learning baselines on the same missingness pattern;
	\item implement and train an LSTM-VAE for sequence reconstruction;
	\item evaluate the methods using masked-point error metrics such as MAE and RMSE;
	\item inspect whether the learned model actually preserves local structure better than simpler alternatives.
\end{itemize}

The central research question is direct: when the same temperature series is corrupted in the same way, does the LSTM-VAE produce a better reconstruction than the baselines, and if so, under what kinds of gaps does that advantage appear most clearly? The answer may not be identical across all missingness patterns. In fact, one useful outcome of the study may be discovering where the deep model does not help as much as expected.

The scope is deliberately narrow. The work concentrates on a single sensor and a single temperature channel so that the comparison remains interpretable and the effect of each imputation method stays visible. That choice leaves multivariate sensor fusion, irregular sampling, and streaming deployment outside the present study. Those extensions are interesting, but they would blur the central comparison and make the chapter harder to read without strengthening the core argument.

Within that boundary, the thesis makes three modest but useful contributions. First, it sets out a fair experimental protocol in which all methods face the same mask and the same ground truth. Second, it collects statistical, machine-learning, and deep-learning imputers in one place so their differences can be compared without shifting assumptions from one implementation to another. Third, it explores an LSTM-VAE design for time-series imputation that is trained and evaluated in a windowed fashion, which may make the method easier to scale than a single full-sequence pass.

The remainder of the thesis is organised as follows. Chapter 2 reviews related work on missing-data mechanisms, time-series imputation, and deep learning approaches for sensor data. Chapter 3 describes the dataset, masking strategy, baseline methods, and the proposed LSTM-VAE pipeline. Chapter 4 presents the experimental results and compares the methods quantitatively and visually. Chapter 5 discusses what the results suggest, where the methods appear to fail, and how the design might be extended. Chapter 6 concludes the thesis and points to directions for future work.
